\documentclass[11pt,a4paper]{article}
\usepackage[T1]{fontenc}
\usepackage{lmodern}
\usepackage[margin=27mm]{geometry}
\usepackage{amsmath,amssymb,amsthm,mathtools}
\usepackage{microtype}
\usepackage{booktabs,array}
\usepackage{enumitem}
\usepackage[hidelinks]{hyperref}
\setlength{\parskip}{5pt plus 1pt}
\setlength{\parindent}{0pt}
\setlength{\emergencystretch}{3em}
\setlist{itemsep=2pt,topsep=3pt}
\raggedbottom
\newtheorem{theorem}{Theorem}[section]
\newtheorem{lemma}[theorem]{Lemma}
\newtheorem{corollary}[theorem]{Corollary}
\newtheorem{proposition}[theorem]{Proposition}
\theoremstyle{definition}
\newtheorem{definition}[theorem]{Definition}
\newtheorem{example}[theorem]{Example}
\theoremstyle{remark}
\newtheorem{remark}[theorem]{Remark}
\newcommand{\Z}{\mathbb{Z}}
\newcommand{\N}{\mathbb{N}}
\newcommand{\Q}{\mathbb{Q}}
\newcommand{\F}{\mathbb{F}}
\newcommand{\ps}[1]{#1[[x]]}
\newcommand{\coeff}[2]{[x^{#1}]\,#2}
\newcommand{\hl}[1]{\textbf{#1}}
\newcommand{\para}[1]{\par\medskip\noindent\textbf{#1}\enspace}
\hypersetup{pdftitle={Every coefficient is 1 modulo 3: a proof of the conjecture in OEIS A396102},pdfauthor={Dev D. Goyal},pdfsubject={OEIS A396102}}
\title{Every coefficient is $1$ modulo $3$: the series with $A(A(A(x)))=(1+x)\,A(A(x))$\\[6pt]\large A proof of the conjecture in OEIS A396102, and of a family of similar statements}
\author{Dev D. Goyal}
\date{October 2026}
\begin{document}
\maketitle
\begin{abstract}
The OEIS entry A396102, contributed by Paul D.~Hanna, concerns the power series
$A(x)=x+x^2-2x^3+10x^4-71x^5+616x^6-\cdots$ determined by the equation $A(A(A(x)))=(1+x)\cdot A(A(x))$, which relates
the second and third iterates of $A$ (the series obtained by substituting $A$ into itself). Hanna conjectured that
every coefficient leaves remainder $1$ on division by $3$. We prove this, as one case of a family of results: for every
$r\ge1$, the solution of $F^{\circ r}=(1+x)\cdot F^{\circ(r-1)}$ agrees modulo $r$ with $x/(1-x)=x+x^2+x^3+\cdots$; the
case $r=3$ is Hanna's conjecture. The key step is to pass to the compositional inverse of the solution, which satisfies
a simpler equation, and to observe that $x/(1+x)$ solves that equation modulo $r$. As by-products, every term of OEIS
A140094 is $1$ modulo $3$, and every term of OEIS A067145 is odd. The main results have been checked by computer and
proved formally in the Lean proof assistant.
\end{abstract}
\begin{center}\begin{minipage}{0.88\linewidth}\small
\textbf{How this paper was produced.} This paper reports the result of an autonomous AI research effort.
ChatGPT (OpenAI), using the GPT-6 model with Pro reasoning and working from prompts by Dev~D.~Goyal, found the proof, wrote its Lean~4 formalization, and wrote a first draft of the paper.
Dev~D.~Goyal chose the problems and asked for checks and revisions, but did not supply proof ideas.
Claude Opus~5.5, used with Claude Code, then checked the work independently: it rebuilt the formal proofs from
source and confirmed that they rely on no unproved assumptions, recomputed the sequence from its definition and tested
every claim against the published data, and read every step of the proofs. It also rewrote this paper for clarity.
Dev~D.~Goyal has verified the results and is responsible for them. The paper has not been peer reviewed.
The section ``About this paper'' at the end gives the details.
\end{minipage}\end{center}

\section{Introduction}

\subsection{The sequence}\label{sec:sequence}
The \emph{On-Line Encyclopedia of Integer Sequences} (OEIS)~\cite{oeisdb} is a large public database of integer
sequences, used by mathematicians and scientists to identify sequences and to record what is known about them. Many
entries also record \emph{conjectures}: patterns that the contributor observed in the data but could not prove.

This paper is about entry A396102~\cite{oeis}, contributed by Paul~D.~Hanna in June 2026. Its name is:
\begin{quote}\small\ttfamily\raggedright G.f. A(x) satisfies A(A(A(x))) = (1+x) * A(A(x)).\par\smallskip\hfill\textrm{--- OEIS A396102, Paul D. Hanna, June 2026}\end{quote}

Here ``G.f.''\ stands for \emph{generating function}: the series $A(x)=\sum_{n\ge1}a(n)\,x^n$ whose coefficients are
the terms of the sequence. The expressions $A(A(x))$ and $A(A(A(x)))$ are \emph{iterates}: $A$ substituted into itself
once and twice (Section~\ref{sec:background}). The factor $(1+x)$ is an ordinary multiplication. So the equation says:
substituting $A$ into itself twice gives the same series as substituting it once and then multiplying by $1+x$.
Together with the normalization $A(x)=x+\cdots$, this determines $A$ completely (Lemma~\ref{lem:aux}). The first
terms, from $n=1$, are 1, 1, $-$2, 10, $-$71, 616, $-$6119, 67210, $-$799307, 10152043, $-$136403576, 1925530303, \dots{}; the listed terms alternate in sign from $n=2$ on.

The entry also records that the compositional inverse $R$ of $A$ (the series with $R(A(x))=x$) satisfies the simpler
equation $R(x)=x/(1+R(R(R(x))))$, and that $-R(-x)$ is the generating function of an older sequence, OEIS
A140094~\cite{oeisR}. It compares $A$ with the series of OEIS A067145~\cite{oeisF2}, which satisfies $F(F(x))=(1+x)F(x)$.

\subsection{The question}
The entry contains the following conjecture.
\begin{quote}\small\ttfamily\raggedright Conjecture: a(n) == 1 (mod 3) for n >= 1.\par\smallskip\hfill\textrm{--- OEIS A396102, Paul D. Hanna, June 2026}\end{quote}

Here \texttt{a(n) == 1 (mod 3)} means that $a(n)$ leaves remainder $1$ on division by $3$. The table shows the
remainders of the first terms modulo $3$, and, for comparison, modulo $9$. Modulo $3$ every remainder is $1$; modulo $9$
the remainders are $1$, $4$ or $7$ (the only possibilities, given the remainder modulo $3$) without an evident pattern.
\begin{center}\small\setlength{\tabcolsep}{3.2pt}
\begin{tabular}{>{\raggedright\arraybackslash}p{6.4em}>{\centering\arraybackslash}p{1.15em}>{\centering\arraybackslash}p{1.15em}>{\centering\arraybackslash}p{1.15em}>{\centering\arraybackslash}p{1.15em}>{\centering\arraybackslash}p{1.15em}>{\centering\arraybackslash}p{1.15em}>{\centering\arraybackslash}p{1.15em}>{\centering\arraybackslash}p{1.15em}>{\centering\arraybackslash}p{1.15em}>{\centering\arraybackslash}p{1.15em}>{\centering\arraybackslash}p{1.15em}>{\centering\arraybackslash}p{1.15em}>{\centering\arraybackslash}p{1.15em}>{\centering\arraybackslash}p{1.15em}>{\centering\arraybackslash}p{1.15em}>{\centering\arraybackslash}p{1.15em}}
\toprule
$n$ & 1 & 2 & 3 & 4 & 5 & 6 & 7 & 8 & 9 & 10 & 11 & 12 & 13 & 14 & 15 & 16 \\ \midrule
$a(n)\bmod 3$ & 1 & 1 & 1 & 1 & 1 & 1 & 1 & 1 & 1 & 1 & 1 & 1 & 1 & 1 & 1 & 1 \\
$a(n)\bmod 9$ & 1 & 1 & 7 & 1 & 1 & 4 & 1 & 7 & 1 & 7 & 1 & 4 & 7 & 7 & 7 & 4 \\
\bottomrule
\end{tabular}
\\[4pt]
\begin{tabular}{>{\raggedright\arraybackslash}p{6.4em}>{\centering\arraybackslash}p{1.15em}>{\centering\arraybackslash}p{1.15em}>{\centering\arraybackslash}p{1.15em}>{\centering\arraybackslash}p{1.15em}>{\centering\arraybackslash}p{1.15em}>{\centering\arraybackslash}p{1.15em}>{\centering\arraybackslash}p{1.15em}>{\centering\arraybackslash}p{1.15em}>{\centering\arraybackslash}p{1.15em}>{\centering\arraybackslash}p{1.15em}>{\centering\arraybackslash}p{1.15em}>{\centering\arraybackslash}p{1.15em}>{\centering\arraybackslash}p{1.15em}>{\centering\arraybackslash}p{1.15em}>{\centering\arraybackslash}p{1.15em}>{\centering\arraybackslash}p{1.15em}}
\toprule
$n$ & 17 & 18 & 19 & 20 & 21 & 22 & 23 & 24 & 25 & 26 & 27 & 28 & 29 & 30 & 31 & 32 \\ \midrule
$a(n)\bmod 3$ & 1 & 1 & 1 & 1 & 1 & 1 & 1 & 1 & 1 & 1 & 1 & 1 & 1 & 1 & 1 & 1 \\
$a(n)\bmod 9$ & 7 & 7 & 1 & 7 & 4 & 7 & 4 & 7 & 1 & 4 & 4 & 4 & 1 & 4 & 4 & 7 \\
\bottomrule
\end{tabular}
\end{center}

\subsection{What this paper proves}
We study the whole family of equations
\begin{equation}\label{eq:family}
F^{\circ r}(x)=(1+x)\cdot F^{\circ(r-1)}(x),\qquad F(x)=x+\cdots\qquad(r\ge1),
\end{equation}
where $F^{\circ r}$ is the $r$-th iterate of $F$. The case $r=3$ is the equation of A396102, and the case $r=2$ is the
equation $F(F(x))=(1+x)F(x)$ of A067145.

\textbf{Main result} (Theorem~\ref{thm:main}). \emph{For every $r\ge1$, equation~\eqref{eq:family} has exactly one
solution $F_r$ with integer coefficients, and modulo $r$}
\[
F_r(x)\equiv\frac{x}{1-x}=x+x^2+x^3+\cdots\pmod r .
\]
\emph{So every coefficient of $F_r$ leaves remainder $1$ on division by $r$. More generally, the coefficient of $x^n$ in
the $k$-th iterate of $F_r$ is congruent to $k^{\,n-1}$ modulo $r$.}

For $r=3$ this proves Hanna's conjecture (Corollary~\ref{cor:conj}). We also deduce that every term of A140094 is
$1$ modulo $3$ and every term of A067145 is odd (Section~\ref{sec:discussion}).

\subsection{What was known}
The equation for the inverse series $R$ recorded in the entry (its formula~(5)) is the starting point of our proof; we
prove it from the definition (Lemma~\ref{lem:aux}). In October 2026, Amol Bhave added to OEIS A196523~\cite{oeisB} a
proof of a congruence for a different family, $A(x)=x+x\cdot A^{\circ l}(x)$, by comparing the solution with
$x/(1-x)$ in the same way. The equation treated here is different: it relates two consecutive iterates, and the
comparison works for the inverse series rather than for $A$ itself. We are not aware of an earlier proof of the
conjecture of A396102.

\subsection{Why it is interesting}
Equations that involve iterates of an unknown series, such as $A(A(A(x)))=(1+x)A(A(x))$, are hard to study: there is
no general method for solving them, and the coefficients of the solutions are usually large and irregular. Simple
remainders are one of the few kinds of structure that can be proved. This example shows that the structure can be
explained by a single simple series, $x/(1-x)$, whose iterates are $x/(1-kx)$. It also shows the value of looking at a
whole family of equations at once, and at the inverse series: the proof becomes a short calculation, and it gives
results for other OEIS entries at no extra cost.

\subsection{How to read this paper}
Section~\ref{sec:background} explains the few notions we need: power series, substitution, compositional inverses,
remainders, and the series $x/(1-cx)$. Section~\ref{sec:results} states the results precisely, and
Section~\ref{sec:proofs} proves them. Section~\ref{sec:checks} describes how the results were checked, by computer and
by a formal proof in Lean. Section~\ref{sec:discussion} discusses consequences for two related OEIS entries and some
open questions.
\section{Background}\label{sec:background}
This section collects the definitions and facts used later. Readers who know formal power series and congruences can skip it and return when needed.
\subsection{Formal power series}
A \emph{formal power series} is an expression
\[
U(x)=u_0+u_1x+u_2x^2+u_3x^3+\cdots ,
\]
determined by its sequence of coefficients $u_0,u_1,u_2,\dots$. No number is ever substituted for $x$, so questions
of convergence do not arise. We write $\coeff{n}{U}=u_n$ for the coefficient of $x^n$ in $U$; two series are equal
when all their coefficients are equal.

Series are added coefficient by coefficient and multiplied like polynomials:
\[
\coeff{n}{(UV)}=u_0v_n+u_1v_{n-1}+\cdots+u_nv_0 .
\]
Each coefficient of a product is a finite sum. For example, $(1-x)(1+x+x^2+x^3+\cdots)=1$, so $1+x+x^2+\cdots$ is
the \emph{multiplicative inverse} of $1-x$, written $\frac{1}{1-x}$. More generally, every series $D$ with constant term
$1$ has a unique multiplicative inverse $\frac1D$, whose coefficients are determined one at a time by $D\cdot\frac1D=1$
using only addition, subtraction and multiplication. Every fraction in this paper has a denominator of this kind, so it
denotes a power series.

We write $U=O(x^d)$ when all coefficients of $U$ below degree $d$ are zero. So
$U-V=O(x^d)$ means that $U$ and $V$ agree up to degree $d-1$. A series of the form $x+u_2x^2+u_3x^3+\cdots$, with
constant term $0$ and coefficient of $x$ equal to $1$, is called \emph{normalized}.

\subsection{Substituting one series into another}
If $V(x)$ has constant term $0$, we can substitute $V$ for $x$ in $U$:
\[
U(V(x))=u_0+u_1V(x)+u_2V(x)^2+u_3V(x)^3+\cdots .
\]
This is called \emph{composition}, also written $U\circ V$. It makes sense because $V(x)^k$ has no terms below $x^k$,
so only the finitely many terms with $k\le n$ affect the coefficient of $x^n$. For example, if $U=x+x^2$ and $V=x^2$, then
$U(V(x))=x^2+x^4$. Composition is associative: $U\circ(V\circ W)=(U\circ V)\circ W$.

Substituting a series into itself repeatedly gives its \emph{iterates}: $U^{\circ2}(x)=U(U(x))$,
$U^{\circ3}(x)=U(U(U(x)))$, and so on, with $U^{\circ0}(x)=x$ and $U^{\circ1}=U$. They must not be confused with
\emph{powers} $U^2=U\cdot U$, $U^3=U\cdot U\cdot U$. (The OEIS sometimes writes $A^k(x)$ for the $k$-th iterate. In
this paper the small circle always means an iterate, and a plain exponent always means a power.)

\subsection{Undoing a substitution}
A normalized series can always be ``undone'' by another normalized series.

\begin{lemma}[Compositional inverse]\label{lem:inverse}
For every normalized series $U$ there is exactly one normalized series $V$ with $U(V(x))=x$, and this $V$ also
satisfies $V(U(x))=x$. If $U$ has integer coefficients, so does $V$.
\end{lemma}

We call $V$ the \emph{compositional inverse} of $U$. The OEIS calls it the \emph{series reversion} of $U$. For example,
the compositional inverse of $x-x^2$ is $C(x)=x+x^2+2x^3+5x^4+14x^5+\cdots$, whose coefficients are the Catalan numbers.
Indeed $C-C^2=x$, which says exactly that substituting $C$ into $x-x^2$ gives $x$.

\begin{proof}
Write $V=x+v_2x^2+v_3x^3+\cdots$ with unknown coefficients. In
$U(V)=V+u_2V^2+u_3V^3+\cdots$, the coefficient of $x^n$ (for $n\ge2$) equals $v_n$ plus terms that involve only
$v_2,\dots,v_{n-1}$. Setting it equal to $0$ therefore determines $v_n$ from the earlier coefficients, using only
additions, subtractions and multiplications. So there is exactly one $V$ with $U\circ V=x$, and it has integer
coefficients if $U$ does. Applying the same construction to $V$ gives a normalized $W$ with $V\circ W=x$. Then
$U=U\circ(V\circ W)=(U\circ V)\circ W=W$, so $V\circ U=V\circ W=x$ as well.
\end{proof}

\subsection{Remainders and congruences}
For an integer $m\ge2$ we write $a\equiv b\pmod m$ (``$a$ is congruent to $b$ modulo $m$'') when $a-b$ is divisible by
$m$, that is, when $a$ and $b$ leave the same remainder after division by $m$. For example, $17\equiv5\pmod 4$, since both leave remainder $1$, and $-3\equiv1\pmod4$. The
remainder of $a$ on division by $m$ is called its \emph{residue} modulo $m$; it is one of $0,1,\dots,m-1$. Saying that
$a$ is odd is the same as saying $a\equiv1\pmod2$.

Addition and multiplication are compatible with congruences, so the residues $0,1,\dots,m-1$, added and multiplied
modulo $m$, form a ring, written $\Z/m\Z$. When $m=p$ is a prime we also write $\F_p$. In $\F_2=\{0,1\}$ we have $1+1=0$,
so addition and subtraction coincide.

For power series with integer coefficients, $U\equiv V\pmod m$ means that \emph{every} coefficient of $U-V$ is
divisible by $m$. Equivalently, the two series become equal when every coefficient is replaced by its residue; we call
this \emph{reducing} the series modulo $m$. The key fact is the following.

\emph{Reduction modulo $m$ is compatible with all the operations above.} To reduce a sum, a product, a substitution, or
an inverse, we may reduce the ingredients first and then compute in $\Z/m\Z$. This holds because every
coefficient of the result is obtained from finitely many coefficients of the ingredients by adding, subtracting and
multiplying, and these operations are compatible with remainders.

\para{Rings.} Several lemmas below hold over every \emph{commutative ring}, that is, every set with an addition,
subtraction and multiplication satisfying the usual rules of arithmetic (for example $\Z$, $\Q$ and $\Z/m\Z$). No
division is used; this matters because in $\Z/m\Z$ division is not always possible (for instance, $2$ has no inverse in
$\Z/4\Z$).

\subsection{The method}
Almost every proof in this paper follows the same three steps.
\begin{enumerate}
\item \emph{Uniqueness.} Show that the defining equation determines the coefficients one at a time, using only
addition, subtraction and multiplication. Then the equation has exactly one solution of the required shape, in every
commutative ring, including $\Z/r\Z$.
\item \emph{Reduction.} Reducing the integer solution modulo $r$ gives a solution of the reduced equation,
because reduction is compatible with all operations.
\item \emph{A good guess.} Find a simple series that satisfies the reduced equation and check that it does. By uniqueness,
this simple series \emph{is} the reduction of the sequence, and its coefficients can simply be read off.
\end{enumerate}
The creative part is step~3: finding a simple series that works. Step~1 sometimes needs care, when the equation seems to
require a division.
\subsection{The series \texorpdfstring{$x/(1-cx)$}{x/(1-cx)}}
For a number $c$ in a commutative ring, put
\[
M_c(x)=\frac{x}{1-cx}=x+cx^2+c^2x^3+c^3x^4+\cdots=\sum_{n\ge1}c^{\,n-1}x^n .
\]
So the coefficient of $x^n$ in $M_c$ is $c^{n-1}$ (with $c^0=1$). For example $M_1=x/(1-x)=x+x^2+x^3+\cdots$ has every
coefficient equal to $1$, $M_0=x$, and $M_{-1}=x/(1+x)=x-x^2+x^3-\cdots$.

\begin{lemma}\label{lem:mobius}
For all $c,d$ in a commutative ring, $M_c\bigl(M_d(x)\bigr)=M_{c+d}(x)$. Consequently the $k$-th iterate of $M_c$ is
$M_{kc}$, and $M_{-c}$ is the compositional inverse of $M_c$.
\end{lemma}

\begin{proof}
Substituting $y=M_d(x)=x/(1-dx)$ into $M_c(y)=y/(1-cy)$ and multiplying the numerator and the denominator by $1-dx$
gives
\[
M_c(M_d(x))=\frac{x}{(1-dx)-cx}=\frac{x}{1-(c+d)x}=M_{c+d}(x).
\]
(All denominators have constant term $1$, so these fractions are power series.) Applying this $k-1$ times gives
$M_c^{\circ k}=M_{kc}$, and $M_c(M_{-c}(x))=M_0(x)=x$.
\end{proof}

\section{Results}\label{sec:results}

\begin{theorem}\label{thm:main}
For every integer $r\ge1$, equation~\eqref{eq:family} has exactly one solution $F_r$ among the normalized power series
with integer coefficients, and
\[
F_r\equiv\frac{x}{1-x}\pmod r .
\]
More generally, for all $k\ge0$ and $n\ge1$, the coefficient of $x^n$ in the $k$-th iterate satisfies
\[
\coeff{n}{F_r^{\circ k}}\equiv k^{\,n-1}\pmod r .
\]
\end{theorem}

\begin{corollary}[Hanna's conjecture]\label{cor:conj}
The series $A$ of OEIS A396102 is $F_3$. Hence $a(n)\equiv1\pmod3$ for every $n\ge1$, and the coefficients of
the $k$-th iterate of $A$ satisfy $\coeff{n}{A^{\circ k}}\equiv k^{\,n-1}\pmod3$.
\end{corollary}

For instance, the second iterate $A(A(x))=x+2x^2-2x^3+11x^4-80x^5+\cdots$, listed in the entry, has coefficients
$\equiv2^{n-1}\pmod3$, that is, $1,2,1,2,1,\dots$; and the third iterate $A(A(A(x)))=x+3x^2+0x^3+9x^4-69x^5+\cdots$ has
all coefficients after the first divisible by $3$.

\section{Proofs}\label{sec:proofs}

\para{The idea in brief.} Equation~\eqref{eq:family} is awkward to work with directly. Its compositional inverse $R$
satisfies the simpler equation $R=x-R\cdot R^{\circ r}$ (Lemma~\ref{lem:aux}), which determines $R$ one coefficient at
a time. Modulo $r$, the series $x/(1+x)$ solves this equation: its $r$-th iterate is $x/(1+rx)$, which is $x$
modulo $r$. So, by uniqueness, $R\equiv x/(1+x)$, and therefore $F_r$, the inverse of $R$, is congruent to $x/(1-x)$.

\begin{lemma}\label{lem:aux}
Let $r\ge1$, and work in any commutative ring.
\begin{enumerate}[label=(\alph*)]
\item There is exactly one power series $R$ with constant term $0$ that satisfies
\begin{equation}\label{eq:aux}
R=x-R\cdot R^{\circ r},
\end{equation}
and it is normalized.
\item A normalized series $F$ satisfies~\eqref{eq:family} if and only if its compositional inverse satisfies~\eqref{eq:aux}.
\end{enumerate}
Consequently~\eqref{eq:family} has exactly one normalized solution, namely the compositional inverse of $R$, and
reducing the integer solution modulo $r$ gives the solution over $\Z/r\Z$.
\end{lemma}

\begin{proof}
(a) We show that the coefficient of $x^n$ in $R\cdot R^{\circ r}$ depends only on the coefficients of $R$ below degree
$n$. Then~\eqref{eq:aux} determines $R$ one coefficient at a time, using only additions, subtractions and
multiplications, so it has exactly one solution. Since $R\cdot R^{\circ r}$ has no terms below $x^2$, the solution
starts with $x$, that is, it is normalized.

Suppose $G$ and $H$ have constant term $0$ and agree below degree $N$, that is, $G-H=O(x^N)$. Then
$G(U)-H(U)=(G-H)(U)=O(x^N)$ for every series $U$ with constant term $0$, and $H(U)-H(V)=O(x^N)$ whenever
$U-V=O(x^N)$, because each $U^k-V^k$ is divisible by $U-V$. Applying these two facts repeatedly shows that all
iterates of $G$ and $H$ agree below degree $N$. Then
\[
G\cdot G^{\circ r}-H\cdot H^{\circ r}=G\,\bigl(G^{\circ r}-H^{\circ r}\bigr)+(G-H)\,H^{\circ r}=O(x^{N+1}),
\]
because in each product one factor is $O(x^N)$ and the other has no constant term. Taking $N=n$ gives the claim.

(b) Let $F$ be normalized with inverse $R$, so that $R(F(x))=x$ and $F(R(x))=x$ (Lemma~\ref{lem:inverse}). Suppose $R$
satisfies~\eqref{eq:aux}. Substitute $F^{\circ r}(x)$ for $x$ in~\eqref{eq:aux}. Since $R\circ F^{\circ r}=F^{\circ(r-1)}$
and $R^{\circ r}\circ F^{\circ r}=x$, this gives
\[
F^{\circ(r-1)}=F^{\circ r}-F^{\circ(r-1)}\cdot x,
\]
which is~\eqref{eq:family}. Conversely, suppose $F$ satisfies~\eqref{eq:family}, and substitute $R^{\circ r}(x)$ for $x$.
Since $F^{\circ r}\circ R^{\circ r}=x$ and $F^{\circ(r-1)}\circ R^{\circ r}=R$, this gives $x=(1+R^{\circ r})\cdot R$,
which is~\eqref{eq:aux}. (Substitution is compatible with sums and products, so we may substitute into both sides of an
identity.)

Finally, by (a) and (b) the normalized solutions of~\eqref{eq:family} are exactly the inverses of the one solution of
\eqref{eq:aux}, so there is exactly one. Reduction modulo $r$ is compatible with all the operations involved, so it maps
the integer solution of~\eqref{eq:aux} to the solution over $\Z/r\Z$, and it maps compositional inverses to
compositional inverses.
\end{proof}

\begin{proof}[Proof of Theorem~\ref{thm:main}]
Work over $\Z/r\Z$ and use the series $M_c=x/(1-cx)$ of Lemma~\ref{lem:mobius}. We check that $M_{-1}=x/(1+x)$
solves~\eqref{eq:aux}. By Lemma~\ref{lem:mobius}, $M_{-1}^{\circ r}=M_{-r}$, and $-r=0$ in $\Z/r\Z$, so
$M_{-1}^{\circ r}=M_0=x$. Therefore
\[
x-M_{-1}\cdot M_{-1}^{\circ r}=x-\frac{x}{1+x}\cdot x=\frac{x(1+x)-x^2}{1+x}=\frac{x}{1+x}=M_{-1}.
\]
By the uniqueness in Lemma~\ref{lem:aux}, the reduction of $R$ modulo $r$ is $M_{-1}$. Its compositional inverse is
$M_1$ (Lemma~\ref{lem:mobius}), so $F_r\equiv M_1=x/(1-x)\pmod r$. Finally, reduction is compatible with substitution,
so $F_r^{\circ k}\equiv M_1^{\circ k}=M_k=\sum_{n\ge1}k^{\,n-1}x^n\pmod r$.
\end{proof}

\begin{proof}[Proof of Corollary~\ref{cor:conj}]
The series $A$ is normalized and satisfies~\eqref{eq:family} with $r=3$, so $A=F_3$ by the uniqueness in
Lemma~\ref{lem:aux}. Apply Theorem~\ref{thm:main} with $r=3$.
\end{proof}

\section{Checking the results}\label{sec:checks}

A mathematical proof can contain mistakes, so the results were checked in two independent ways.

\para{Numerical check.}
We solved the defining equation coefficient by coefficient, with exact integer arithmetic. This reproduces all 23
terms printed in the OEIS entry. The congruence $a(n)\equiv1\pmod3$ holds for every $n<70$, and so does
$\coeff{n}{A^{\circ k}}\equiv k^{\,n-1}\pmod3$ for $0\le k\le8$. For the family, we checked that every coefficient of
$F_r$ is $1$ modulo $r$, for $2\le r\le7$ and $n<40$.
The computation was done with the PARI/GP computer algebra system, using code that was written separately from the
formal proofs (file \texttt{verification/independent\_checks.gp} in the accompanying repository~\cite{repo}).
A computation can only test finitely many cases, so it is a check on the definitions and on the indexing, not a proof.

\para{Formal proof in Lean.}
Lean~\cite{lean} is a \emph{proof assistant}: a program in which mathematical statements and their proofs are written
in a precise formal language, and every step is checked mechanically down to the basic rules of logic. All results of
this paper are proved in Lean~4 (version 4.19.0) on top of Mathlib~\cite{mathlib}, Lean's library of formalized
mathematics. The statements are collected in the file \texttt{lean/Entries/A396102.lean} of the repository,
in the namespace \texttt{OEIS.A396102}. In Lean the family is indexed so that \texttt{family\_source\_existence r} is the case $r+1$ of~\eqref{eq:family}. The proofs are in the file \texttt{lean/Families/InverseIterate.lean}. The consequences for OEIS A140094 and A067145 in Section~\ref{sec:discussion} are not part of the Lean formalization.
The table lists where each result of this paper is proved.
\begin{center}\small\renewcommand{\arraystretch}{1.25}
\begin{tabular}{>{\raggedright\arraybackslash}p{0.52\linewidth}>{\raggedright\arraybackslash}p{0.4\linewidth}}
\toprule
Statement in this paper & Lean theorem(s)\\
\midrule
$A$ is normalized, satisfies $A(A(A(x)))=(1+x)A(A(x))$, and is the only such series & \texttt{normalized}, \texttt{source\_\allowbreak{}equation}, \texttt{source\_\allowbreak{}unique}, \texttt{source\_\allowbreak{}existence\_\allowbreak{}uniqueness} \\
Corollary~\ref{cor:conj}: $a(n)\equiv1\pmod3$, and the iterates & \texttt{coefficient\_\allowbreak{}congruence}, \texttt{iterate\_\allowbreak{}congruence} \\
Lemma~\ref{lem:aux}: the family~\eqref{eq:family} and the equation~\eqref{eq:aux} for the inverse & \texttt{family\_\allowbreak{}source\_\allowbreak{}existence}, \texttt{family\_\allowbreak{}source\_\allowbreak{}equation}, \texttt{auxiliary\_\allowbreak{}source\_\allowbreak{}equation}, \texttt{inverse\_\allowbreak{}right}, \texttt{inverse\_\allowbreak{}left} \\
Theorem~\ref{thm:main} for every $r$ & \texttt{family\_\allowbreak{}coefficient\_\allowbreak{}congruence}, \texttt{family\_\allowbreak{}iterate\_\allowbreak{}congruence} \\
\bottomrule
\end{tabular}
\end{center}
The Lean project was rebuilt from source for this paper, and Lean was asked to list the basic assumptions (axioms)
that each of these theorems ultimately depends on. Every one depends only on the three standard axioms of Lean's logic
(\texttt{propext}, \texttt{Classical.choice}, \texttt{Quot.sound}); no step is left unproved. A formal proof certifies
exactly the statement that Lean checks. Whether that statement says what the OEIS entry means is a separate question,
so the Lean definitions were also compared, side by side, with the definition in the OEIS entry
(Section~\ref{sec:sequence}).

\section{Discussion}\label{sec:discussion}

\para{Scope.} Theorem~\ref{thm:main} determines every coefficient of every member of the family modulo $r$, and every
coefficient of every iterate. It says nothing about higher powers of $r$: modulo $9$ the coefficients of A396102 show no
evident pattern (see the table in Section~\ref{sec:sequence}). It also says nothing about the signs. The terms of
A396102 alternate in sign from $n=2$ on for every $n<60$ that we computed; we have not proved this.

\para{Two related entries.} The theorem has consequences for two older OEIS entries, which we checked against their
published terms.
\begin{itemize}
\item \emph{OEIS A140094}~\cite{oeisR}, by Paul~D.~Hanna (2008), has a generating function $S$ with
$S(x)=x/(1-S(S(S(x))))$. Let $R$ be the inverse of $A$ as above, and put $S(x)=-R(-x)$, as recorded in A396102. From
$R(-y)=-S(y)$ one gets $R^{\circ 3}(-x)=-S^{\circ3}(x)$, and substituting $-x$ for $x$ in $R=x-R\cdot R^{\circ3}$ gives
$S=x+S\cdot S^{\circ3}$, which is the same as $S=x/(1-S^{\circ3})$. (By the argument of Lemma~\ref{lem:aux}(a), this
equation has only one solution with constant term $0$, so $S$ is indeed the generating function of A140094.) Since $R\equiv x/(1+x)\pmod3$, we get
$S(x)\equiv x/(1-x)\pmod3$. So \emph{every term of A140094 is $1$ modulo $3$.} The 23 terms printed in the entry
agree with this. The same argument shows that, for every $r$, the solution of $S=x/(1-S^{\circ r})$ is congruent to
$x/(1-x)$ modulo $r$.
\item \emph{OEIS A067145}~\cite{oeisF2}, by Christian~G.~Bower (2002), is the sequence that ``shifts left under
reversion'': its generating function $F$ satisfies $R(x)=F(x)/x-1$, where $R$ is its inverse. Substituting $F(x)$ for
$x$ turns this into $F(F(x))=(1+x)F(x)$, as noted in that entry; this is~\eqref{eq:family} with $r=2$. So $F=F_2$,
and by Theorem~\ref{thm:main} \emph{every term of A067145 is odd.} The 23 terms printed in the entry agree with this.
\end{itemize}

\para{Why $x/(1-x)$?} The series $M_1=x/(1-x)$ is special because its iterates are so simple: $M_1^{\circ k}=x/(1-kx)$.
Modulo $r$, the $r$-th iterate is $x$ again, so in any equation in which iterates appear only through their number
modulo $r$, the series $x/(1-x)$ or $x/(1+x)$ is a natural candidate. The same observation underlies Bhave's proof
for A196523~\cite{oeisB}. It would be interesting to know which other iterative equations in the OEIS can be treated
this way.

\section*{About this paper}\label{sec:about}
\para{How the work was done.}
This paper is one of 29 papers on conjectures from the OEIS that came out of an autonomous AI research effort directed
by Dev~D.~Goyal~\cite{repo}. The roles were as follows.
\begin{itemize}
\item \emph{Dev~D.~Goyal} started the project, chose which problems to pursue, asked for proofs, formal verification and
revisions, and decided what to publish. Dev~D.~Goyal did not supply proof strategies or intermediate arguments.
\item \emph{ChatGPT} (OpenAI), using the GPT-6 model with Pro reasoning, worked autonomously from those prompts. It
searched the OEIS for suitable conjectures, found the proofs, wrote the Lean~4 formalization, and wrote a first draft
of this paper.
\item \emph{Claude Opus~5.5}, used with Claude Code, checked the work independently and then rewrote this paper. It
rebuilt the Lean project from source and confirmed that the theorems listed in Section~\ref{sec:checks} rely only on
Lean's standard axioms; compared the Lean definitions with the OEIS definition; recomputed the sequence with separately
written PARI/GP code, compared it with the terms published in the OEIS, and tested every claim of the paper on many
terms; and read every step of the proofs. It then rewrote the text to make it self-contained and readable. The
theorems and the proof strategies are ChatGPT's. In rewriting, Claude Opus~5.5 expanded the explanations and, in
places, wrote out an argument in more detail or by a more elementary route. The Discussion section was added by Claude Opus~5.5;
the further observations in it were checked by computer and are not part of the Lean formalization.
\end{itemize}
Dev~D.~Goyal has verified the results and is responsible for this paper.

\para{Review status.} This paper has not been peer reviewed, and it has not been submitted to a journal or to the OEIS.

\para{Priority.} An OEIS entry that labels a statement as a conjecture shows only that its contributor did not know
a proof; it does not show that no proof existed. During the research only limited searches for earlier work were
made. No claim of priority is made.

\para{Corrections.} Please report any error in the issue tracker of the repository~\cite{repo}, where the Lean sources and the numerical checks are also
available.

\section*{Acknowledgements}
The sequence and the conjecture are due to Paul~D.~Hanna~\cite{oeis}, who also recorded the equation for the inverse
series and the link to A140094. The comparison with $x/(1-x)$ is also the idea of Amol~Bhave's proof in
A196523~\cite{oeisB}. OEIS A067145 was contributed by Christian~G.~Bower.

\begin{thebibliography}{99}\small
\raggedright
\bibitem{oeis} P.~D.~Hanna, entry A396102 (G.f.\ $A(x)$ satisfies $A(A(A(x)))=(1+x)A(A(x))$), \emph{The On-Line Encyclopedia of Integer Sequences}, June 2026. \url{https://oeis.org/A396102}. Text as of October 2026.
\bibitem{oeisR} P.~D.~Hanna, entry A140094 (G.f.\ satisfies $A(x)=x/(1-A(A(A(x))))$), \emph{The On-Line Encyclopedia of Integer Sequences}, 2008. \url{https://oeis.org/A140094}. Text as of October 2026.
\bibitem{oeisF2} C.~G.~Bower, entry A067145 (Shifts left under reversion), \emph{The On-Line Encyclopedia of Integer Sequences}, 2002. \url{https://oeis.org/A067145}. Text as of October 2026.
\bibitem{oeisB} P.~D.~Hanna, entry A196523 (G.f.\ $A(x)$ satisfies $A(x)=x+x\,A(A(A(A(x))))$), \emph{The On-Line Encyclopedia of Integer Sequences}, 2011. \url{https://oeis.org/A196523}. Comment by A.~Bhave, 13 October 2026. Text as of October 2026.
\bibitem{oeisdb} N.~J.~A.~Sloane and The OEIS Foundation Inc., \emph{The On-Line Encyclopedia of Integer Sequences}. \url{https://oeis.org}.
\bibitem{stanley} R.~P.~Stanley, \emph{Enumerative Combinatorics}, Volume~2, Cambridge University Press, 1999. (Chapter~5 treats formal power series, composition and compositional inverses.)
\bibitem{niven} I.~Niven, Formal power series, \emph{American Mathematical Monthly} 76 (1969), 871--889.
\bibitem{lean} L.~de~Moura and S.~Ullrich, The Lean~4 theorem prover and programming language, in \emph{Automated Deduction -- CADE 28}, Lecture Notes in Computer Science 12699, Springer, 2021, 625--635.
\bibitem{mathlib} The mathlib Community, The Lean mathematical library, in \emph{Proceedings of the 9th ACM SIGPLAN International Conference on Certified Programs and Proofs (CPP 2020)}, 367--381. Version used: commit \texttt{c44e0c8ee63c}.
\bibitem{repo} D.~D.~Goyal, \emph{Proofs and corrections of OEIS conjectures, with Lean~4 formalizations}, repository, 2026. \url{https://github.com/dvgyl/oeis-a396843-proof}.
\end{thebibliography}
\end{document}
